\section{Experimental Methodology}

The System Architecture chapter established that BENEFITS is a running system, not a specification: the EDC connector extensions, the consent packages, and the audit and privacy pipelines all execute against real infrastructure and pass a two-language conformance-vector suite, C1, that proves the TypeScript edge-side evaluator and the Java connector-side evaluator agree on every conformance case. C1 is a correctness baseline, not a performance or robustness claim: it establishes that consent is evaluated identically wherever it is evaluated, but it says nothing about what that evaluation costs, how the architecture behaves when its assumptions are stressed, or where its guarantees degrade under realistic operating conditions. Six mechanism-level experiments build on that baseline to answer exactly those questions. Each experiment targets one open question the architecture's design left unsettled at specification time, each is driven by a seeded, deterministic harness whose full parameter grid and repeated trials produce reproducible results from a committed seed, and each writes its output to a committed CSV under \texttt{experiments/results/}, so every number reported in Chapter 5 traces back to a specific, rerunnable sweep rather than a one-off measurement. This chapter describes what each experiment measures, how its parameter space is swept, and what statistics are computed over the results. It does not report or characterize the results themselves, which are reserved for Chapter 5.

\subsection{E1: Signature-Verification Cost at the Edge}

The architecture's participant-side enforcement (S1) is a \texttt{validate\_doc\_update} function installed on each participant's own CouchDB database, evaluated inside CouchDB's own JavaScript query-server engine rather than on-device. SIDS 2026 left open whether this design could be extended with cryptographic assurance, specifically whether a pure-JavaScript Ed25519 signature verification routine, running inside that same constrained query-server engine, is fast enough to be viable at any realistic batch size or fails outright. E1 answers this by measuring real HTTP writes against a real CouchDB instance under three arms: a baseline with no validator installed (isolating CouchDB's own request and storage overhead), the existing structural validator unchanged, and a third, experiment-only design document that performs the same structural check and additionally verifies an Ed25519 signature over each record. The signature-verifying arm exists only inside the experiment harness. It is never installed by any production provisioning path.

Each arm is swept across three dimensions: batch size (10, 50, 100, 500, 1000, 5000, 10000 documents, log-spaced), field count per record (1, 5, 20), and grant-set size (0, 10, or 50 decoy grants added ahead of each field's real grant, chosen specifically because consent-core's grant lookup scans linearly and returns on first match, so decoy count, not field count, is what determines how far that scan actually has to walk). The full cross product is 63 configurations per arm, each repeated across 3 trials to give the results a basis for reporting variance rather than a single point per cell, run against a throttled container fixed at one CPU core and 1 GB of memory to approximate the class of low-end handset the target population is expected to own, with the host CPU model recorded alongside every row so the throttle quota is interpretable rather than assumed reproducible across machines. Wall-clock time is measured client-side around each write batch.

\subsection{E2: Revocation Propagation Under Network Partition}

The device-fleet phase of the reference implementation already proved, at the mechanism level, that a participant who revokes consent while offline with records already queued produces zero leaked records once connectivity resumes, and that conflicting concurrent writes converge under the project's "user preference first" resolution strategy. E2 is the formal, swept version of that same mechanism: it asks how revocation propagation latency and time-to-consistency vary with connectivity, not merely whether the guarantee holds at one hand-built connectivity pattern.

The experiment simulates a fleet of 10 devices per trial, split into three fixed roles: three revocation devices, each of which queues an ordinary record and then revokes consent for it while still offline, before any sync has occurred. Three conflict devices, each of which writes a document under a known identifier and then has a rival remote revision injected against it. And four baseline devices that contribute ordinary records with no revocation or conflict, giving the trial's leaked-record and consistency checks real successful writes to check against. Each device runs its own duty cycle governing when it comes online, and the one swept dimension is that duty cycle's online probability: 0.05, 0.1, 0.3, 0.5, and 0.9, spanning near-permanent offline to near-permanent online. Thirty independent fleet trials run at each of the five sweep points, for 150 trials in total, each trial running for up to 200 simulated ticks. After each device completes its tracked action, the harness checks the shared CouchDB database directly for any leaked record matching a revoked participant and purpose, and checks that every conflict document's conflict list has been fully resolved. Either check failing increments a defect counter rather than being folded into a graceful statistic. For the timing behavior, the harness records the median and interquartile range of propagation ticks (for revocation devices) and consistency ticks (for conflict devices) pooled across all 90 data points collected per sweep point, alongside the leaked-record and unresolved-conflict counts, which the design treats as pass/fail correctness gates rather than as distributions to characterize.

\subsection{E3: Audit-Chain Verification Cost at Scale}

The audit-chain mechanism's boolean correctness, that a mutated leaf, root, or proof step fails verification, and that a reordered or altered epoch is caught by chain verification, was already proven at the unit level during an earlier implementation phase. What that phase deliberately deferred was the question SIDS 2026 left open: epoch length trades proof size against sealing latency, and the reference implementation needed to know what that tradeoff actually costs on the participant's own device as the audit structure scales. E3 supplies that measurement without touching the correctness logic itself. It adds timing and size instrumentation around calls that already exist.

Two dimensions are swept. Epoch size (10, 50, 100, 500, 1000, 5000 entries, log-spaced) governs proof size and inclusion-proof verification time: for each epoch size, the harness builds one Merkle tree and samples inclusion proofs at 10 evenly spaced leaf positions (or every leaf, below 10), recording the number of proof steps and the resulting proof size in digest bytes. Because a single inclusion-proof verification is sub-millisecond even at the largest epoch size, each fixture is timed across 10 repetitions after an untimed warmup pass, and the median of those repetitions is recorded, following the same repeated-trial convention E1 uses to make a sub-millisecond measurement reportable at all. Chain length (5, 20, 100, 500 epochs) governs chain verification time separately, since verifying the chain only recomputes and compares epoch-to-epoch hash links and its cost does not depend on epoch size at all. Sweeping it against epoch size would produce no informational content, so it is measured against chain length instead, using the same warmup-then-median-of-ten-repetitions protocol. All fixture generation is seeded so results reproduce from a committed seed, and inclusion-proof verification runs under the same throttled, low-end device profile as E1, since that is the code path a participant's own device actually executes. This experiment follows most directly from \cite{crosby2009}, who established proof size and verification cost as the metric of interest for tree-structured tamper-evident logs and measured both as a function of log scale, and it sits alongside \cite{melara2015} and \cite{hu2021}, whose Merkle-based transparency logs report comparable proof-size- and latency-versus-scale curves under their own scale parameters. E3 departs from all three by sweeping epoch size specifically rather than overall log or user-base scale, since epoch size is the parameter this architecture's own sealing design exposes as a tunable tradeoff.

\subsection{E4: Privacy-Budget-Versus-Utility Tradeoff}

SIDS 2026 raised, as an open governance challenge rather than a settled design, the concern that differential privacy could render aggregate statistics analytically meaningless at the population scale of an individual service site. E4 answers this directly by driving the privacy pipeline's existing, unmodified k-anonymity, Laplace-noise, and budget-accountant classes with synthetic query workloads and measuring accuracy as a function of privacy budget, population size, and the k-anonymity threshold, rather than modifying any of the production classes under test.

A seeded population generator produces synthetic contributions with three quasi-identifier attributes (service type, channel, urgency level), under two distribution modes: uniform, and a skewed mode in which channel preference is fixed at 61 percent paper contact, matching the one workshop statistic from the parent study that maps directly onto a swept attribute. The remaining attribute skews are declared assumptions rather than workshop-sourced figures, stated as such. Eight queries are run per sweep point in a fixed order against one shared budget accountant per synthetic site: three grouped counts (over service type, channel, and urgency level, chosen for their differing category cardinalities), two proportions (the largest service-type group's share, and a fixed rare category's share), and three cross-tabulations of increasing granularity. Running all eight against one shared accountant, rather than giving each query its own fresh budget, is deliberate: it is the only way to exercise the accountant's real, atomic, shared-budget behavior, which is exactly the property under test.

The sweep varies epsilon (the per-query privacy parameter: 0.1, 0.27, 0.31, 0.37, 0.5, 1.0, 1.5, and 2.0, with the three added values chosen specifically to fall inside the epsilon range where the number of queries a fixed total per-site budget of 2.0 can still afford actually changes, since a coarser five-point grid was found to sit entirely on one side of that transition and produce a degenerate result), site population (30, 100, 200, 300), and the k-anonymity threshold (3, 5, 7, 10), crossed with the two population-skew modes, for 256 total sweep points. Because the computation is deterministic given a seed, with no wall-clock variability to average out, each point runs once rather than across repeated trials. The median and interquartile range the parent design calls for are instead computed across the pooled per-group relative-error values the eight-query workload produces at that point. Three statistics are reported per point: the median and interquartile range of relative error across the full pooled set, with a fixed sentinel value of 1.0 assigned to any group that was suppressed by k-anonymity or any query refused for exhausted budget. The suppression/refusal rate itself, isolating how much of the workload received no answer at all. And a surviving-only median and interquartile range computed over only the subset of data points that were neither suppressed nor refused, isolating the accuracy of the answers that were actually returned from how often an answer was returned at all. This experimental shape, accuracy characterized as a function of privacy budget under a shared, enforced accountant, follows the precedent set by \cite{abadi2016}, who report the accuracy cost of a tracked privacy budget directly, and \cite{hay2016}, whose DPBench methodology establishes accuracy-versus-epsilon curves under a shared benchmark protocol as the standard empirical framing for this class of evaluation. \cite{rogers2021} supplies the closest production analogue, a deployed budget-management service enforcing a strict, shared differential-privacy budget on live queries at scale, structurally the same kind of accountant this experiment measures at a much smaller, per-site scale.

\subsection{G1: Quorum Availability Under the Built Policy and the Veto Policy SIDS 2026 Specified}

The steward-quorum mechanism requires a threshold number of stewards to co-sign before a data release proceeds. The built system's production decision function reads only the count of signing stewards. It never consults recorded vetoes, even though SIDS 2026's own text specifies a veto as part of the release gate, not merely a quorum: "any single steward may veto, in which case the release request is returned to the requester with a written justification" \cite{katrakazas2026}. G1 exists to answer what that gap leaves open: what does enforcing the veto guarantee SIDS 2026 specified, and the built system currently does not provide, cost in release availability as stewards go offline, act slowly, or object.

Two arms run on an identical, shared simulated timeline of steward actions. Arm A calls the real, unmodified quorum-decision class directly against a real, unmodified steward-authorization store as the timeline unfolds, the built system's own veto-blind policy. Arm B is a second decision function that exists only inside the experiment harness, explicitly documented as non-production, which adds exactly the check SIDS 2026 specified and the built system omits: a permanent block the moment any single steward vetoes. At each discrete tick, every steward who has not yet acted may act with a fixed per-tick probability, and, conditional on acting, may sign or veto according to a fixed veto share. A run terminates when both arms reach a determinate outcome, quorum reached for Arm A, quorum reached, permanently vetoed, or starved for Arm B, or at a 200-tick horizon chosen so that even the least favorable swept parameters leave the probability of an undetermined outcome negligible.

The sweep spans steward pool size (5, 10, 20), quorum size expressed as a fraction of the pool (0.3, 0.5, 0.7), per-tick action probability (0.1, 0.3, 0.5, 0.7, 0.9), and veto share among acting stewards (0.0, 0.1, 0.3, with 0.0 serving as a sanity check where Arm B degenerates to Arm A's own behavior), for 135 sweep points, each run across 200 independent seeded repetitions to characterize the distribution of a stochastic per-tick process rather than a single deterministic run. Per point, the harness reports each arm's reached rate, and, for Arm A, its starved rate, for Arm B, its reached, vetoed, and starved rates, which are mutually exclusive and sum to one, so that Arm B's failures can be decomposed into how much comes from vetoes actually landing versus genuine starvation rather than reported as one undifferentiated failure rate. Time-to-quorum is reported as the median and interquartile range of ticks, computed only over the runs that actually reached quorum in that arm, following the same surviving-only convention E4 established for its own accuracy statistic, so that an arm's speed and its reliability are reported as separate, uncombined statistics. This design mirrors the quorum-size-versus-performance evaluation shape established by \cite{kokoriskogias2020}, whose CALYPSO system measures transaction latency as trustee-committee size scales, the closest direct precedent in the literature for characterizing a threshold-committee's behavior as a tunable tradeoff between committee size and system outcome, here reliability rather than latency, under a two-arm comparison between the built, veto-blind policy and the veto-enforcing policy SIDS 2026's own text specifies.

\subsection{G2: Consent Heterogeneity Versus Aggregate Utility}

The architecture's consent model gates writes but, before this experiment, gated nothing at query time: an already-written contribution had no mechanism preventing its inclusion in an aggregate query result regardless of what purposes or recipient classes its owner had actually consented to. G2 closes that gap as a real, additive production capability, then uses it to answer a question the parent design's meet operation raised but never measured: whether realistic, heterogeneous per-participant consent choices make community aggregates too thin to be useful, as distinct from the noise the privacy mechanism itself introduces.

A population generator, built on the same seeded generator E4 uses, additionally assigns each participant a consent-grant profile. With probability \texttt{1 - h}, where h is the heterogeneity parameter under test, a participant is given a broad grant set covering every attribute the query workload needs. With probability h, one attribute, chosen at random, is left without a valid grant, excluding that participant from any query that groups by that attribute while leaving them included in queries that do not need it. At h = 0, every participant consents to everything the workload asks, reproducing the implicit assumption every earlier experiment made without stating it. At h = 1, every participant is missing coverage for exactly one field.

The sweep varies epsilon (0.25, 0.5, 1.0, 2.0, with 0.25 added because it is the value at which the full eight-query workload clears the same fixed total per-site budget of 2.0 that E4 uses, needed to make the heterogeneity-driven exclusion signal observable rather than swamped by budget refusal), population (100, 300), and heterogeneity h (0.0, 0.1, 0.2, 0.3, 0.5, 0.7, 1.0, finely spaced near the low end where the initial slope of the utility-versus-heterogeneity relationship matters most), with k fixed at 5 and population skew fixed at uniform, since neither is this experiment's focus, for 56 total sweep points, each deterministic given its seed and run once. At each point, every query in the workload is evaluated against two separate ground truths: one over the full generated population, and one restricted to participants who actually hold a valid, matching grant. Two relative-error statistics are computed against these respectively, both from the same noised release: coverage error, comparing that release against the consenting population as reference and so isolating the privacy mechanism's own noise and suppression behavior, and exclusion error, comparing the same release against the full-population ground truth and so capturing noise and exclusion bias together. Both are reported as pooled medians and interquartile ranges across the workload, alongside a surviving-only pair of each statistic computed over only the data points that were neither suppressed nor refused, and the excluded fraction, a direct, epsilon-independent measurement of how much of the generated population lacked coverage for at least one query's attributes at a given heterogeneity level. Because exclusion error nests coverage error's noise component plus an additional exclusion-bias term, it is the gap between exclusion error and coverage error, not exclusion error in isolation, that isolates the effect of heterogeneous consent specifically. That gap is the design choice that lets this experiment answer SIDS 2026's underlying question cleanly: whether a thin aggregate, if one appears, is a consequence of differential-privacy noise, of heterogeneous consent excluding participants, or of both, rather than reporting a single conflated utility number.

\#\# Closing

Taken together, these six experiments cover the architecture's principal points of designed uncertainty: cryptographic enforcement cost at the participant boundary, correctness of consent revocation under realistic connectivity failure, the audit chain's cost as it scales, the accuracy the privacy pipeline can deliver under a shared, enforced budget, the availability behavior of the steward-quorum release gate under two competing policy designs, and the effect of realistic, non-uniform consent choices on aggregate query utility. Each experiment's design was chosen to follow an established empirical evaluation shape from the relevant systems and privacy literature rather than an ad hoc protocol invented for this paper, and each is fully specified by a seed, a fixed parameter grid, and a committed set of statistics, so that the results reported in the next chapter are reproducible from the harnesses and CSVs this methodology describes.
